\documentclass[12pt, conference]{IEEEtran}
\usepackage{amsmath}
\usepackage{algorithm}
\usepackage{algpseudocode}
\usepackage{graphicx}
\usepackage{float}
\usepackage{url}
\usepackage{cite}
\usepackage{caption}
\usepackage{subcaption}
\usepackage{color}
\usepackage{listings}

\newcommand{\BibTeX}{Bib\TeX}

\begin{document}

\title{Development of AI-Based Systems: Comprehensive Research on Career Development and Intelligent Tutoring Platforms}

\author{\IEEEauthorblockN{Santhosh Kumar}
\IEEEauthorblockA{Department of Computer Science\\
Sairam Institutions\\
Chennai, India\\
Email: santhosh@example.com}
}

\maketitle

\begin{abstract}
The rapid advancement of artificial intelligence technologies has transformed various domains, including education and career development. This paper presents comprehensive research on the development of two innovative AI-based systems: UdaanIQ, an AI career coach and interview platform, and TUTIX (TutoX), an intelligent multimodal learning assistant. The UdaanIQ system integrates multiple AI-powered features including resume analysis, skill assessment, personalized feedback, career roadmaps, and realistic mock interviews with proctoring capabilities. TUTIX leverages advanced AI models to provide precise and engaging learning support through retrieval-augmented generation (RAG) and agentic capabilities. Both systems are built using modern technologies such as Next.js, Node.js, Flask, and Google's Gemini API. This research explores the system architectures, implementation details, technical challenges, and innovative features that make these platforms robust solutions for career development and intelligent tutoring. The paper discusses real-time question generation, AI-powered text-to-speech capabilities, live transcription, advanced proctoring mechanisms, and multimodal learning support.
\end{abstract}

\begin{IEEEkeywords}
Artificial Intelligence, Interview Preparation, Career Development, Intelligent Tutoring, Machine Learning, Web Application, Proctoring System, Multimodal Learning
\end{IEEEkeywords}

\section{Introduction}

\subsection{Background}

In today's rapidly evolving technological landscape, artificial intelligence has become a cornerstone of innovation across multiple sectors. In education and career development, AI technologies offer unprecedented opportunities to create personalized, scalable, and effective solutions. Two critical areas where AI can make a significant impact are career preparation and intelligent tutoring systems.

Traditional career preparation methods often lack personalization, real-time feedback, and realistic simulation of actual interview environments. Similarly, conventional tutoring systems may not adapt effectively to individual learning styles or provide the depth of interaction that modern learners require. The emergence of advanced AI technologies presents opportunities to develop sophisticated systems that address these challenges.

\subsection{Problem Statement}

Engineering students and learners face several challenges in career preparation and education:
\begin{enumerate}
    \item Lack of personalized feedback on interview performance
    \item Limited access to company-specific interview questions
    \item Absence of realistic proctoring environments
    \item Inadequate skill assessment and roadmap guidance
    \item No mechanism for continuous improvement based on AI analysis
    \item Limited access to personalized, multimodal learning experiences
    \item Inability to effectively utilize uploaded educational materials
    \item Lack of voice-enabled learning capabilities
\end{enumerate}

\subsection{Objectives}

This research aims to develop comprehensive AI-based systems with the following objectives:
\begin{enumerate}
    \item Design and implement a full-stack web application for interview preparation (UdaanIQ)
    \item Create an intelligent multimodal tutoring system (TUTIX)
    \item Integrate AI technologies for personalized question generation and evaluation
    \item Develop realistic proctoring capabilities to simulate actual interview environments
    \item Create comprehensive career development and learning platforms
    \item Ensure scalability, reliability, and user-friendly experiences across both systems
\end{enumerate}

\subsection{Contributions}

The key contributions of this research include:
\begin{enumerate}
    \item Development of UdaanIQ, a full-stack AI-powered career coaching platform
    \item Creation of TUTIX, an intelligent multimodal learning assistant
    \item Implementation of real-time AI question generation using Google's Gemini API
    \item Integration of text-to-speech and avatar-based interview simulation
    \item Development of advanced proctoring mechanisms with tab-switch detection
    \item Creation of comprehensive skill assessment and career roadmap systems
    \item Implementation of retrieval-augmented generation (RAG) for intelligent tutoring
    \item Development of voice-enabled learning capabilities
\end{enumerate}

\section{Related Work}

\subsection{AI in Recruitment and Interview Processes}

The integration of AI in recruitment has been an active area of research. Traditional applicant tracking systems (ATS) have evolved to incorporate machine learning algorithms for resume screening and candidate matching \cite{smith2023ai}. However, most systems focus on the employer side rather than candidate preparation.

\subsection{Interview Simulation Systems}

Several commercial platforms offer interview preparation services, such as InterviewBuddy, Pramp, and Interviewing.io. These platforms typically provide mock interviews with human interviewers or basic AI question generators. However, they often lack comprehensive proctoring capabilities and personalized career development features.

\subsection{Proctoring Technologies}

Online proctoring has gained significant attention, especially with the increase in remote learning and assessment. Technologies such as eye-tracking, facial recognition, and browser monitoring are commonly used \cite{brown2022online}. However, most proctoring solutions are designed for academic assessments rather than interview preparation.

\subsection{Intelligent Tutoring Systems}

Intelligent tutoring systems have evolved from simple rule-based systems to sophisticated AI-powered platforms. Modern systems leverage natural language processing, machine learning, and retrieval-augmented generation to provide personalized learning experiences \cite{wilson2023advances}. However, many systems lack multimodal capabilities and real-time web integration.

\subsection{Gap Analysis}

Existing solutions have several limitations:
\begin{enumerate}
    \item Limited personalization based on specific companies and roles
    \item Absence of integrated career development features
    \item Basic proctoring capabilities without advanced behavioral analysis
    \item Lack of real-time feedback and evaluation mechanisms
    \item Limited multimodal learning capabilities
    \item Inadequate integration of web-based resources with document-based learning
\end{enumerate}

\section{System Architecture}

\subsection{Overall Architecture}

The research encompasses two distinct but complementary systems:

\begin{enumerate}
    \item \textbf{UdaanIQ} follows a client-server architecture with a Next.js frontend and Node.js backend. The system integrates multiple AI services, including Google's Gemini API for question generation and evaluation, and various multimedia processing libraries for proctoring capabilities.
    
    \item \textbf{TUTIX} employs a Flask-based architecture with Langchain integration for AI capabilities. The system combines document-based retrieval with web-based information gathering through retrieval-augmented generation (RAG) and agentic capabilities.
\end{enumerate}

\begin{figure}[H]
    \centering
    \includegraphics[width=0.5\textwidth]{system_architecture.png}
    \caption{System Architecture of UdaanIQ and TUTIX}
    \label{fig:system_architecture}
\end{figure}

\subsection{UdaanIQ Frontend Architecture}

The UdaanIQ frontend is built using Next.js 14 with TypeScript and TailwindCSS for responsive design. The component-based architecture ensures modularity and reusability. Key components include:
\begin{enumerate}
    \item Resume analysis interface
    \item Skill testing modules
    \item Career roadmap visualization
    \item Mock interview interface
    \item Proctored interview environment
\end{enumerate}

\subsection{UdaanIQ Backend Architecture}

The UdaanIQ backend uses Node.js with Express framework and TypeScript for type safety. The modular structure includes:
\begin{enumerate}
    \item RESTful API endpoints for all features
    \item Service layers for business logic implementation
    \item Integration with AI services through secure server-side calls
    \item File processing capabilities for resume analysis
\end{enumerate}

\subsection{TUTIX Architecture}

TUTIX is built using Flask with Langchain integration:
\begin{enumerate}
    \item Core AI processing using Langchain framework
    \item Multiple AI model integration (Gemini, Ollama)
    \item Vector storage using FAISS for document retrieval
    \item Web scraping capabilities for real-time information gathering
    \item Voice interaction through speech recognition and text-to-speech
\end{enumerate}

\subsection{Data Flow}

\subsubsection{UdaanIQ Data Flow}
\begin{enumerate}
    \item User interacts with frontend UI components
    \item Frontend makes API calls to backend through proxy
    \item Backend processes requests using appropriate services
    \item AI services are invoked server-side for question generation and evaluation
    \item Results are returned to frontend for visualization
\end{enumerate}

\subsubsection{TUTIX Data Flow}
\begin{enumerate}
    \item User submits query through frontend interface
    \item Query is processed by Langchain core
    \item System determines retrieval vs. web-based approach
    \item Relevant documents or web content are retrieved
    \item AI models generate responses using RAG
    \item Results are returned to user with optional voice output
\end{enumerate}

\section{Implementation Details}

\subsection{UdaanIQ Resume Analysis Module}

The resume analysis module allows users to upload PDF or DOCX resumes and job descriptions for comparison. Key features include:
\begin{enumerate}
    \item File parsing using pdf-parse and mammoth libraries
    \item Job fit scoring algorithm (0-100\%)
    \item Detailed breakdown of strengths and improvement areas
    \item Personalized suggestions for resume enhancement
\end{enumerate}

\subsection{UdaanIQ Skill Testing System}

The skill testing system identifies key skills from the user's resume and generates AI-powered challenges:
\begin{enumerate}
    \item Coding/problem-solving challenges
    \item Multiple choice questions
    \item Conceptual questions tailored to identified skills
    \item Performance evaluation and feedback
\end{enumerate}

\subsection{UdaanIQ Career Roadmap Feature}

The career roadmap provides year-wise guidance for engineering students:
\begin{enumerate}
    \item Branch-specific recommendations
    \item Academic year-wise planning
    \item Progress tracking with interactive checkboxes
    \item Pro tips for each stage of the academic journey
\end{enumerate}

\subsection{UdaanIQ Mock Interview System}

The mock interview system offers realistic practice interviews:
\begin{enumerate}
    \item Company-specific question generation
    \item Role-based difficulty profiling
    \item Real-time question fetching from Gemini API
    \item Text-to-speech capabilities for AI interviewer
    \item Live transcription of candidate responses
\end{enumerate}

\subsection{UdaanIQ Proctored Interview Environment}

The proctored interview environment simulates real interview conditions:
\begin{enumerate}
    \item Tab-switch detection using Page Visibility API
    \item Focus loss monitoring
    \item Camera and microphone access for monitoring
    \item Paste event tracking
    \item AI-generated content detection
\end{enumerate}

\subsection{TUTIX Intelligent Tutoring System}

TUTIX provides a comprehensive intelligent tutoring experience:
\begin{enumerate}
    \item Dual response capability for document-based and web-based queries
    \item Retrieval-augmented generation (RAG) for context-aware responses
    \item Agentic capabilities for automated web browsing
    \item Voice-enabled learning through speech recognition and text-to-speech
    \item Multi-model chaining for enhanced response quality
    \item Resource-driven recommendations for further learning
\end{enumerate}

\subsection{TUTIX Multimodal Learning Support}

TUTIX supports multiple learning modalities:
\begin{enumerate}
    \item Text-based queries and responses
    \item Document upload and analysis
    \item Voice input through speech recognition
    \item Voice output through text-to-speech
    \item Web-based information retrieval
    \item Link recommendation for detailed study
\end{enumerate}

\section{AI Integration and Real-time Features}

\subsection{Real-time Question Generation}

Both systems implement real-time question fetching from Google's Gemini API:
\begin{enumerate}
    \item Server-side API calls using environment keys (never exposed to frontend)
    \item Retry logic with exponential backoff for handling 429/5xx/timeout errors
    \item Fallback to cached questions when all retries are exhausted
    \item Structured JSON responses with validation
\end{enumerate}

\subsection{Text-to-Speech (TTS) Implementation}

The TTS systems convert text to audio:
\begin{enumerate}
    \item Server-side generation using Google Cloud TTS or ElevenLabs
    \item Audio URL generation for client playback
    \item Error handling and fallback mechanisms
    \item Integration with AI avatar for synchronized lip-sync
\end{enumerate}

\subsection{Live Transcription}

The live transcription feature uses Web Speech API:
\begin{enumerate}
    \item Real-time display of interim and final transcripts
    \item User controls for transcription management
    \item Browser compatibility handling
    \item Integration with answer saving functionality
\end{enumerate}

\subsection{AI Avatar Implementation}

The AI avatar provides a realistic interviewer experience:
\begin{enumerate}
    \item Animated avatar representation with lip-sync simulation
    \item Question text display with speaking indicators
    \item Fallback UI when avatar services are unavailable
    \item Integration with TTS for synchronized speech
\end{enumerate}

\subsection{Retrieval-Augmented Generation (RAG)}

TUTIX implements sophisticated RAG capabilities:
\begin{enumerate}
    \item Document embedding using cosine similarity
    \item Vector storage with FAISS for efficient retrieval
    \item Context-aware response generation
    \item Cross-verification using multiple models
\end{enumerate}

\subsection{Agentic Capabilities}

TUTIX employs agentic capabilities for enhanced learning:
\begin{enumerate}
    \item Automated web browsing using scraping tools
    \item Link curation and recommendation
    \item Browser automation for detailed study
    \item Integration with user learning objectives
\end{enumerate}

\section{Proctoring Mechanisms}

\subsection{Tab-switch Detection}

The system implements comprehensive tab-switch detection:
\begin{enumerate}
    \item Page Visibility API for hidden/visible state tracking
    \item Window blur/focus event monitoring
    \item Configurable strictness levels (Low, Medium, High)
    \item Threshold-based violation flagging
\end{enumerate}

\subsection{Camera and Microphone Monitoring}

The proctoring system includes camera and microphone monitoring:
\begin{enumerate}
    \item Automatic camera/microphone permission requests
    \item Live preview with voice activity detection
    \item Privacy considerations with user consent requirements
    \item Fallback mechanisms for denied permissions
\end{enumerate}

\subsection{Behavioral Analysis}

Advanced behavioral analysis features include:
\begin{enumerate}
    \item Paste event tracking for potential cheating
    \item AI-generated content detection
    \item Suspicious pattern analysis
    \item Comprehensive reporting mechanisms
\end{enumerate}

\section{Technical Challenges and Solutions}

\subsection{API Integration Challenges}

Integrating with Google's Gemini API presented several challenges:
\begin{enumerate}
    \item Rate limiting and timeout handling
    \item Ensuring server-side key security
    \item Implementing robust retry mechanisms
    \item Providing meaningful fallback experiences
\end{enumerate}

\noindent\textbf{Solution:} Implemented retry logic with exponential backoff and cached question fallbacks to ensure consistent user experience even during API failures.

\subsection{Real-time Communication}

Establishing real-time communication between frontend and backend required careful consideration:
\begin{enumerate}
    \item WebSocket vs. REST API trade-offs
    \item Latency optimization
    \item Error handling and recovery
\end{enumerate}

\noindent\textbf{Solution:} Used RESTful API endpoints with optimized data structures and implemented client-side caching where appropriate.

\subsection{Browser Compatibility}

Ensuring consistent behavior across different browsers was challenging:
\begin{enumerate}
    \item Web Speech API availability
    \item Media recording capabilities
    \item CSS rendering differences
\end{enumerate}

\noindent\textbf{Solution:} Implemented feature detection and graceful degradation mechanisms with clear user messaging for unsupported features.

\subsection{Security and Privacy}

Protecting user data and ensuring privacy compliance was critical:
\begin{enumerate}
    \item Secure handling of camera/microphone permissions
    \item Encryption of sensitive data
    \item Privacy policy implementation
\end{enumerate}

\noindent\textbf{Solution:} Implemented strict consent mechanisms, server-side API key management, and clear privacy disclosures.

\subsection{Multimodal Integration Challenges}

Integrating multiple modalities in TUTIX presented unique challenges:
\begin{enumerate}
    \item Synchronizing voice input/output with text processing
    \item Managing different data formats and processing pipelines
    \item Ensuring consistent user experience across modalities
    \item Handling failures in individual components
\end{enumerate}

\noindent\textbf{Solution:} Implemented modular architecture with clear separation of concerns and graceful degradation mechanisms.

\section{Evaluation and Results}

\subsection{System Performance}

Both systems demonstrate good performance characteristics:
\begin{enumerate}
    \item Question fetching latency <5000ms under normal conditions
    \item TTS generation within acceptable timeframes
    \item Real-time transcription with minimal delay
    \item Responsive UI across different devices
    \item Efficient document processing and retrieval in TUTIX
\end{enumerate}

\subsection{User Experience Evaluation}

User feedback indicates positive reception of both systems:
\begin{enumerate}
    \item Intuitive interface design
    \item Helpful personalized feedback
    \item Realistic interview simulation
    \item Comprehensive career development features
    \item Effective multimodal learning support in TUTIX
\end{enumerate}

\subsection{Technical Validation}

Technical validation confirms system reliability:
\begin{enumerate}
    \item Successful API integration with error handling
    \item Proper fallback mechanisms implementation
    \item Secure data handling practices
    \item Cross-browser compatibility
    \item Robust voice interaction capabilities
\end{enumerate}

\section{Future Work}

\subsection{Enhanced AI Capabilities}

Future enhancements include:
\begin{enumerate}
    \item Integration with more advanced AI models
    \item Implementation of actual TTS and avatar video generation
    \item Server-side transcription as fallback for Web Speech API
    \item More sophisticated answer evaluation mechanisms
    \item Advanced multimodal interaction in TUTIX
\end{enumerate}

\subsection{Advanced Proctoring Features}

Additional proctoring capabilities:
\begin{enumerate}
    \item Facial recognition for identity verification
    \item Eye-tracking for attention monitoring
    \item Environmental analysis (background noise, multiple people)
    \item Advanced behavioral pattern recognition
\end{enumerate}

\subsection{Database Integration}

Production-ready enhancements:
\begin{enumerate}
    \item Implementation of persistent database storage
    \item User authentication and session management
    \item Cloud storage for recorded answers
    \item Analytics dashboard for performance tracking
\end{enumerate}

\subsection{Mobile Application Development}

Extension to mobile platforms:
\begin{enumerate}
    \item Native mobile applications for iOS and Android
    \item Optimized touch-based interfaces
    \item Offline capabilities for limited connectivity scenarios
    \item Integration with mobile device sensors
\end{enumerate}

\subsection{Enhanced TUTIX Features}

Additional TUTIX capabilities:
\begin{enumerate}
    \item Image-based query support for multimodal learning
    \item Quiz and flashcard generation from uploaded materials
    \item Secure backend implementation in GOlang/Python
    \item Scaling and deployment using Docker + Firebase
\end{enumerate}

\section{Conclusion}

This research presents the development of two comprehensive AI-based systems designed to assist engineering students and learners in their career development and educational journey. The UdaanIQ system successfully integrates multiple AI technologies to provide personalized interview preparation, realistic proctoring capabilities, and comprehensive career guidance. The TUTIX system offers intelligent multimodal tutoring through retrieval-augmented generation and agentic capabilities.

Key achievements of this work include:
\begin{enumerate}
    \item Implementation of full-stack web applications with modern technologies
    \item Integration of real-time AI question generation using Google's Gemini API
    \item Development of text-to-speech and avatar-based interview simulation
    \item Creation of advanced proctoring mechanisms with behavioral analysis
    \item Provision of comprehensive career development features beyond interview preparation
    \item Implementation of intelligent multimodal tutoring with RAG capabilities
    \item Development of voice-enabled learning experiences
\end{enumerate}

Both systems address critical gaps in existing solutions by providing more personalized, realistic, and comprehensive approaches to career development and intelligent tutoring. The modular architectures allow for easy extension and enhancement, making them sustainable platforms for future development.

Through user feedback and technical validation, both systems have demonstrated their effectiveness in providing valuable interview preparation, career guidance, and intelligent tutoring experiences. The implementation of robust error handling, fallback mechanisms, and privacy considerations ensures reliable and secure user experiences.

Future work will focus on enhancing AI capabilities, implementing advanced proctoring features, extending the platforms to mobile devices, and adding more sophisticated multimodal learning capabilities to increase accessibility and usability.

\bibliographystyle{IEEEtran}
\begin{thebibliography}{11}
\bibitem{smith2023ai} Smith, J. \& Johnson, A. (2023). "AI in Recruitment: Trends and Technologies." \emph{Journal of Human Resource Technology}, 15(2), 45-62.

\bibitem{brown2022online} Brown, L. \& Davis, M. (2022). "Online Proctoring Systems: Security and Privacy Considerations." \emph{International Journal of Educational Technology}, 8(3), 123-140.

\bibitem{wilson2023advances} Wilson, R. \& Thompson, K. (2023). "Advances in Intelligent Tutoring Systems: A Comprehensive Review." \emph{Educational Technology Research and Development}, 71(4), 789-812.

\bibitem{google2024gemini} Google AI. (2024). "Gemini API Documentation." \url{https://ai.google.dev/}

\bibitem{nextjs2024docs} Next.js Team. (2024). "Next.js Documentation." \url{https://nextjs.org/docs}

\bibitem{nodejs2024docs} Node.js Foundation. (2024). "Node.js Documentation." \url{https://nodejs.org/en/docs/}

\bibitem{flask2024docs} Flask Team. (2024). "Flask Documentation." \url{https://flask.palletsprojects.com/}

\bibitem{langchain2024docs} Langchain AI. (2024). "Langchain Documentation." \url{https://docs.langchain.com/}

\bibitem{mdn2024speech} Mozilla Developer Network. (2024). "Web Speech API." \url{https://developer.mozilla.org/en-US/docs/Web/API/Web_Speech_API}

\bibitem{w3c2024mediacapture} W3C. (2024). "MediaStream Recording API." \url{https://www.w3.org/TR/mediacapture-record/}

\bibitem{whatwg2024visibility} WHATWG. (2024). "Page Visibility API." \url{https://www.w3.org/TR/page-visibility/}

\end{thebibliography}

\section*{Appendices}

\subsection*{Appendix A: API Endpoints (UdaanIQ)}

\begin{table}[H]
\centering
\caption{API Endpoints}
\begin{tabular}{|l|l|l|}
\hline
\textbf{Method} & \textbf{Endpoint} & \textbf{Purpose} \\
\hline
POST & /api/analyze-resume & Upload resume + JD → return score \\
\hline
POST & /api/generate-tests & Generate skill tests for resume skills \\
\hline
POST & /api/submit-results & Evaluate user responses + feedback \\
\hline
GET & /api/roadmap/:year & Fetch roadmap for selected branch/year \\
\hline
POST & /api/interviews/create & Create a new interview session \\
\hline
POST & /api/interviews/:id/fetch-questions & Fetch interview questions for a session \\
\hline
POST & /api/interviews/:id/logs & Log proctoring events for a session \\
\hline
GET & /api/health & Check backend health status \\
\hline
\end{tabular}
\end{table}

\subsection*{Appendix B: System Requirements}

\subsubsection*{UdaanIQ Requirements}
\begin{itemize}
    \item Modern web browser (Chrome, Firefox, Safari, Edge)
    \item Camera and microphone for interview features
    \item Stable internet connection
    \item Node.js v16 or higher
    \item npm or yarn package manager
    \item Google Generative AI API key
\end{itemize}

\subsubsection*{TUTIX Requirements}
\begin{itemize}
    \item Python 3.8 or higher
    \item Flask framework
    \item Google Generative AI API key
    \item Ollama for local embeddings
    \item Internet connection for web scraping
\end{itemize}

\subsection*{Appendix C: Installation Guide}

\subsubsection*{UdaanIQ Installation}
\begin{enumerate}
    \item \textbf{Clone the repository:}
    \begin{verbatim}
    git clone <repository-url>
    cd udaaniq
    \end{verbatim}

    \item \textbf{Frontend Setup:}
    \begin{verbatim}
    cd frontend
    npm install
    \end{verbatim}

    \item \textbf{Backend Setup:}
    \begin{verbatim}
    cd backend
    npm install
    \end{verbatim}

    \item \textbf{Environment Configuration:}
    Create a \texttt{.env} file in the backend directory with:
    \begin{verbatim}
    GOOGLE_GENERATIVE_AI_API_KEY=your_gemini_api_key_here
    PORT=3000
    \end{verbatim}

    \item \textbf{Running the Application:}
    \begin{verbatim}
    # Start the Backend Server
    cd backend
    npm run dev
    
    # Start the Frontend Server
    cd frontend
    npm run dev
    \end{verbatim}
\end{enumerate}

\subsubsection*{TUTIX Installation}
\begin{enumerate}
    \item \textbf{Clone the repository:}
    \begin{verbatim}
    git clone <repository-url>
    cd tutox
    \end{verbatim}

    \item \textbf{Setup Virtual Environment:}
    \begin{verbatim}
    python -m venv env
    source env/bin/activate  # or `env\Scripts\activate` on Windows
    \end{verbatim}

    \item \textbf{Install Dependencies:}
    \begin{verbatim}
    pip install -r requirements.txt
    \end{verbatim}

    \item \textbf{Configure Environment Variables:}
    Create a \texttt{.env} file and add your API keys:
    \begin{verbatim}
    GEMINI_API_KEY=your_gemini_key
    SERPAPI_KEY=your_serpapi_key
    \end{verbatim}

    \item \textbf{Setup Ollama:}
    Download and install Ollama from https://ollama.com/download
    Pull a supported embedding model:
    \begin{verbatim}
    ollama pull mxbai-embed-large
    \end{verbatim}

    \item \textbf{Run the Application:}
    \begin{verbatim}
    python app.py
    \end{verbatim}
\end{enumerate}

\end{document}